\mychapter{0}{Glossaire}

Ce glossaire présente les principaux termes et concepts utilisés tout au long de ce mémoire.

\vspace{1em}

\begin{description}
    \item[Revue de code (Code Review)] Processus systématique d'examen du code source par des pairs développeurs dans le but de détecter les défauts, d'améliorer la qualité du code et de partager les connaissances au sein d'une équipe.
    
    \item[Pull Request (PR) / Merge Request (MR)] Mécanisme permettant à un développeur de proposer des modifications de code à un projet. Les autres membres de l'équipe peuvent alors examiner, commenter et approuver ces modifications avant leur intégration.
    
    \item[LLM (Large Language Model)] Modèle de langage de grande taille basé sur l'apprentissage profond, capable de comprendre et de générer du texte en langage naturel. Exemples : GPT, Claude, LLaMA.
    
    \item[DevOps] Approche combinant le développement logiciel (Dev) et les opérations informatiques (Ops) visant à raccourcir le cycle de développement et à fournir des livraisons continues de haute qualité.
    
    \item[API (Application Programming Interface)] Interface de programmation permettant à différentes applications de communiquer entre elles. GitHub et GitLab fournissent des APIs pour accéder aux données des dépôts.
    
    \item[Bottleneck (Goulot d'étranglement)] Point de congestion dans un processus qui ralentit l'ensemble du flux de travail. Dans le contexte DevOps, une revue de code inefficace peut devenir un bottleneck.
    
    \item[GitLab] Plateforme web de gestion de versions basée sur Git, offrant des fonctionnalités de revue de code, d'intégration continue et de déploiement continu.
    
    \item[GitHub] Plateforme d'hébergement de code source utilisant Git, proposant des outils de collaboration, de revue de code et de gestion de projets.
    
    \item[Commit] Enregistrement d'une modification dans l'historique d'un dépôt Git, incluant les changements apportés au code ainsi qu'un message descriptif.
    
    \item[Dataset] Ensemble de données collectées et organisées dans un format structuré, utilisé pour l'analyse et la recherche.
    
    \item[Ingénierie Logicielle Empirique] Approche de recherche en génie logiciel basée sur l'observation, l'expérimentation et l'analyse de données réelles provenant de projets logiciels.
    
    \item[Open Source] Logiciel dont le code source est accessible publiquement, permettant à quiconque de l'examiner, le modifier et le distribuer.
    
    \item[Pipeline DevOps] Ensemble automatisé de processus permettant de passer du code source au déploiement en production, incluant les tests, la revue de code, et l'intégration continue.
    
    \item[Pagination] Technique de division de grandes quantités de données en pages plus petites lors de requêtes API, permettant de gérer efficacement les limitations de transfert de données.
    
    \item[Authentification] Processus de vérification de l'identité d'un utilisateur ou d'une application, souvent nécessaire pour accéder aux APIs des plateformes comme GitHub ou GitLab.
    
    \item[RevMine] Outil développé dans le cadre de ce projet, conçu pour faciliter la collecte et l'analyse des données de revue de code à partir de plateformes comme GitHub et GitLab.
    
\end{description}

\clearpage